\section{Scope, standard of proof, and snapshot}
\label{sec:scope}

\subsection{The target is an arithmetic derivation}

The phrase ``formalizable in \(\ipra\)'' is used here in a strict sense.  The
final mathematical assertion must be a first-order formula over natural-number
codes.  Its functions are fixed primitive-recursive symbols, its logic is
intuitionistic, and every use of induction has a quantifier-free induction
formula.  A completed Lean endpoint therefore has the shape
\[
   \Gamma\;\vdash[\texttt{PRAAx}]\;\varphi,
\]
not merely an ambient theorem about Lean's \texttt{Nat}, \texttt{Rat}, or
\texttt{Real}.  The latter can validate an algorithm or an algebraic identity,
but it does not by itself construct a derivation in the embedded theory.

The existing \code{ConciseIPRA} development supplies both sides of this
boundary.  Its ordinary presentation has the types \code{PR} and \code{PTm},
the axiom family \code{PRAAx}, and the derived rule \code{PRA.qfInd}.  Its
native scalar calculus has iterator terms, equation formulas, the two binders
\code{Q} and \code{D}, and derivations \code{NPrf}.  The theorems
\code{pra_iff_native} and \code{native_iff_pra} faithfully transport an
already completed arithmetic derivation.  They are bridges, not automatic
mathematics compilers.

\subsection{Repository snapshot and audit scope}

The public repository was inspected at main-branch commit
\begin{center}
\small\ttfamily
7d492fbfc5fee3b2b9dec0b2623b0b82effe05f2.
\end{center}
The proof-packet index at that snapshot records forty-eight delivered packets
in algebra, analysis, analytic and Diophantine number theory, combinatorics,
geometry, probability, and topology.  The index also warns that packets were
filed as delivered: placement did not rerun their code or compile their Lean
files.  Some contain complete finite arguments; some contain uncompiled
blueprints, explicit \code{sorry} obligations, uninstantiated analytic
contracts, or deliberately isolated impossibility results.

This paper is therefore not a blanket certification of the packets.  It uses
them as a large empirical sample from which to extract common proof shapes and
common failure modes.  Statements about a particular packet are limited to the
paths listed in \cref{app:sources}.  The positive theorems proved below are
framework theorems: finite-trace reflection, regular-name calculus, robust
branching, finite-net search, and certificate transformations.  They do not
supply the missing deep analytic or arithmetic estimates in a packet whose
report explicitly leaves them open.

\subsection{Three notions that must not be conflated}

There are three useful but distinct levels of completion.

\begin{definition}[Semantic implementation]
A semantic implementation is executable code or an ambient Lean theorem that
computes the intended finite object and proves a proposition about its ordinary
meaning.
\end{definition}

\begin{definition}[Arithmetic proof architecture]
An arithmetic proof architecture specifies natural-number codes, fixed
primitive-recursive functions, exact formulas, induction invariants, bounds,
and a dependency-ordered route to an embedded derivation.
\end{definition}

\begin{definition}[Embedded completion]
An embedded completion is an actual derivation in \code{PRAAx}, optionally
transported to \code{NPrf}, with an axiom audit of the final Lean declaration.
\end{definition}

A packet may be strong at the first two levels while incomplete at the third.
The framework is designed to make the transition from the second to the third
mechanical enough that the remaining gap is visible and localized.

\subsection{What ``general'' can mean in a first-order theory}

A textbook theorem often begins ``for every continuous function'' or ``for
every real random variable.''  Literal higher-type quantification is not
available in first-order \(\ipra\).  There are three legitimate replacements.

\begin{enumerate}
\item A \emph{fixed-evaluator schema}: choose a primitive-recursive evaluator
      \(a(\mathbf p,k)\) externally, add its regularity or domain assumptions,
      and prove the resulting first-order formula uniformly in the numerical
      parameters \(\mathbf p\).
\item A \emph{finite-data theorem}: make a rational table, finite law,
      polynomial, graph, triangulation, or certificate the numerical input.
      This is a single uniform first-order theorem.
\item A \emph{trace-relative theorem}: input a finite computation or proof
      trace and prove that acceptance by a primitive-recursive verifier implies
      the desired arithmetic conclusion.  Trace production may be specialized,
      while verification remains uniform.
\end{enumerate}

The last form is central.  It gives generic infrastructure without postulating
a total universal evaluator for all primitive-recursive programs, which is
impossible by diagonalization (\cref{thm:no-universal-pr}).

\subsection{Contribution and limits}

The main contribution is a reusable normal form for positive proofs:
\[
\boxed{\begin{gathered}
\text{fixed numerical input and precision}\;\longrightarrow\;\text{bounded trace}\\
\longrightarrow\;\text{Boolean invariant}\;\longrightarrow\;\text{checked observation}
\end{gathered}}
\]
The exact or asymptotic theorem is recovered only after the finite observation
has been proved for an arbitrary free precision parameter.  This order is the
logical discipline that prevents a hidden use of completeness, countable
choice, or induction on a quantified Cauchy assertion.

The framework resolves many \emph{formulation and compilation} difficulties.
It does not turn ineffective constants into effective ones, produce rational
representatives from abstract Selmer data, or uniformly choose an axis at the
identity rotation.  In those cases it identifies a minimal positive theorem
and isolates the missing producer or the constructive obstruction.
